\SourceLeaf{228}%
Sohn eingeschlagen hat, und die Resultate, zu denen er geführt ist, kommen fast durchgehend mit meinen eigenen, zum Teil schon seit 30—35 Jahren angestellten Meditationen überein. In der Tat bin ich dadurch auf das Äußerste überrascht. Mein Vorsatz war, von meiner eigenen Arbeit, von der übrigens bis jetzt wenig zu Papier gebracht war, bei meinen Lebzeiten gar nichts bekannt werden zu lassen. Die meisten Menschen haben gar nicht den rechten Sinn für das, worauf es dabei ankommt, und ich habe nur wenige Menschen gefunden, die das, was ich ihnen mitteilte, mit besonderem Interesse aufnahmen. Um das zu können, muss man erst recht lebendig gefühlt haben, was eigentlich fehlt, und darüber sind die meisten Menschen ganz unklar. Dagegen war meine Absicht, mit der Zeit alles so zu Papier zu bringen, dass es wenigstens mit mir dereinst nicht unterginge.

Sehr bin ich also überrascht, dass diese Bemühung mir nun erspart werden kann und höchst erfreulich ist es mir, dass gerade der Sohn meines alten Freundes es ist, der mir auf eine so merkwürdige Art zuvorgekommen ist.

Sehr prägnant und abkürzend finde ich die Bezeichnungen. Doch glaube ich, dass es gut sein wird, für manche Hauptbegriffe nicht bloß Zeichen oder Buchstaben, sondern bestimmte Namen festzusetzen, und ich habe bereits vor langer Zeit an einige solcher Namen gedacht. Solange man die Sache nur in unmittelbarer Anschauung durchdenkt, braucht man keine Namen oder Zeichen. Die werden erst nötig, wenn man sich mit anderen verständigen will. So könnte z.~B. die Fläche, die Dein Sohn $F$ nennt, eine Parasphäre, die Linie $L$ ein Paracykel genannt werden. Es ist im Grunde Kugelfläche oder Kreislinie von unendlichem Radius. Hypercykel könnte der Komplexus aller Punkte heißen, die von einer Geraden, mit der sie in einer Ebene liegen, gleiche Distanz haben, ebenso Hypersphäre. Doch das sind alles nur unbedeutende Nebensachen. Die Hauptsache ist der Stoff, nicht die Form.

In manchem Teil der Untersuchung habe ich etwas andere Wege eingeschlagen. Als ein Spezimen füge ich einen rein geometrischen Beweis (in den Hauptzügen) von dem Lehrsatz bei, dass die Differenz der Summe der Winkel eines Dreiecks von $180^{0}$ dem Flächeninhalt des Dreiecks proportional ist.

\Needspace{6\baselineskip}
I. Der Komplexus dreier Geraden $ab$, $cd$, $ef$, die so beschaffen sind, dass $ab\TripleParallel dc$, $cd\TripleParallel fe$, $ef\TripleParallel ba$, bildet eine Figur, die ich $T$ nenne. Es lässt
